\section{Computational Study}\label{sec:study52}
\subsection{Design and independent validation}
The study separates exact structural solvability, additive certification, and numerical optimization. Its protocol fixes all inputs, seeds, tolerances, and comparisons before the full run. It discloses the small development pilots rather than treating them as unseen tests. The original six difficult models are retained at two requested accuracies. Twelve new strict-prefix cases vary history count, catalog size, budget, and promise; six cases vary catalog eligibility under otherwise matched primitives; fourteen cases test within-gap perturbations, threshold crossings, free refinement, and dominated splitting commands. All 44 requests are reported.

Every resource-path interval is checked independently in original policy space. Exact enumeration over books, using a separate antecedent fixed-book allocator, supplies references for all 32 resource-grid cases. The twelve large common-prefix cases instead use the exact pooling theorem and a separate primal-dual and complete-path checker; exhaustive enumeration is not claimed at that scale. A further 36 group-subdivision certificates are checked. This gives 80 independently checked certificates across the principal study and its baseline comparisons.

Regression tests compare the selected-boundary identity with original fixed-book optimization on 1,583 books from 100 seeded models. All 100 corresponding global intervals contain their exhaustive optima. Two thousand concave-kernel convolution cases compare brute force, the optimizer's divide-and-conquer search, and the checker's totally monotone search. Twenty combinatorial tests exercise the dyadic count, including the ratio-three example. Twelve deliberately corrupted certificates are rejected. These finite tests supplement the proofs; they do not substitute for them.

All inputs are synthetic. The computational record includes the exact model vectors, original policies, solver statuses, complete certificate tables, environment, and content hashes. Timing records distinguish optimization from checking. Optimization includes payoff preprocessing, path search, recovery, and table construction; serialization and compression are excluded and disclosed separately. Publication reruns retain the same protocol and generate the displayed tables directly from their execution records. Local development measurements remain separately labeled.

\subsection{Difficult models and competitive comparisons}
Table~\ref{tab:stress52} shows the six models at tolerance 0.001. Every resource-path interval meets the request; the largest width is \FineMaxGap. Optimization takes \FineMinTime--\FineMaxTime\ seconds in the recorded environment. Every returned lower policy equals the exhaustive reference. This latter fact does not turn a positive-width resource certificate into a zero-width global proof: exact attainment is known here from the separate exhaustive comparison.

Absolute accuracy can mask material relative uncertainty. The largest reported relative width is \FineMaxRelative, on a small positive reference value. We therefore report the reference, absolute width, and relative width together, and retain width divided by the primitive scale in the machine-readable record. No economic importance is assigned to a numerical tolerance without application-specific calibration.

Table~\ref{tab:baseline52} compares class-mean subdivision at node allowances 1, 7, 31, and 127, root-price refinement at three settings, and individual-box subdivision. At 127 nodes with the default price setting, \GroupNOneComplete\ of six requests attain 0.001. The original historical record had ten unresolved fine-tolerance requests among its two node budgets; every such request remains archived unchanged. Increasing the root-price set helps some cases but can alter the adaptive cover and worsen another finite-budget interval. We report that outcome rather than selecting the best setting after observing the answer.

These subdivision comparisons use the same machine but fixed node allowances, not matched time and memory ceilings. The direct mixed-integer comparison below shares a measured optimization-time allowance with the resource method. We do not claim a complete equal-memory ranking, and we have not added node-local price generation or an exhaustive repaired-net timing study. The stronger variable-eligibility algorithm removes the need to rely on those heuristics for the new polynomial guarantee.

\subsection{Nontrivial eligibility and catalog refinement}
Table~\ref{tab:strict52} reports exact common-eligibility instances with 16 through 1,024 histories. All ceilings lie strictly inside the same interior catalog gap, differ across histories, and exclude a nonempty suffix. Thus the result is not restricted to the all-ceilings-equal-one geometry. Each case retains nonuniform weights, heterogeneous caps and opening charges, and zero as well as positive quadratic curvatures. Promises include interior and upper-bound obligations. All twelve cases close with zero rational width at the root.

Table~\ref{tab:sweep52} varies the full-catalog class count from one to twelve while holding the other primitives fixed. The generic resource method is run at each count, including the one-class case, to isolate its behavior. The grid dimension does not grow with that class count; edge capacities and exact arithmetic still change, so identical runtime is neither expected nor asserted. The two-class returned policy is below its exact reference by $83/143360$; the other five attain their references. This distinguishes the additive guarantee from exact attainment. The study remains modest in catalog size and does not establish large-instance superiority for variable eligibility.

The refinement experiment inserts four commands into a four-point catalog. With opening charge three, each insertion is provably dominated, yet the full-catalog class count rises from one to five. The winning two-command book, its single selected eligibility class, and its exact value remain unchanged. With free insertions, the first command improves the value and the next three split classes without further improvement; the winning book has only two selected classes. Within-gap ceiling perturbations preserve the value exactly. The near-boundary pair changes full-catalog eligibility but retains its winning book; the analytical example in Section~\ref{sec:robust52} separately demonstrates that a true value jump is possible. The companion reports all fourteen cases, not only those with a changed optimum.

\subsection{Certificate cost and mixed-integer diagnostics}
Table~\ref{tab:cert52} reports certificate economics on the fine-tolerance models. Compressed sizes range from \FineMinPacked\ to \FineMaxPacked\ decimal megabytes; the largest recorded denominator has \FineMaxDenBits\ bits. The checker reconstructs the arc values and checks all anchor tables, so certification carries a real storage and computation cost. A small terminal alphabet is not a bound on the target table, encoder logic, probability precision, or audit storage.

Table~\ref{tab:mip52} exposes the numerical envelope results rather than summarizing only their completion status. The implementation uses SciPy's mixed-integer interface \citep{Virtanen2020} on a direct support-indicator probability/service formulation, independently of the resource-path reduction. All twelve envelope solves finish successfully within their assigned allowances. They are much faster than the fine resource-grid construction on these small instances; the largest paired runtime is \MIPMaxSeconds\ seconds. Exhaustive fixed-book enumeration is also competitive at this scale. The case for the resource method is its proved additive complexity and exact original-space certificates, not a universal speed ranking.

The numerical upper and lower values agree to floating-point precision on these cases, including one tiny negative raw width from rounding. We preserve that diagnostic rather than relabeling it an exact proof. Envelope approximation allowances, solver relative gaps, model residuals, branch-and-bound nodes, and both objectives remain in the record. None of these floating-point statuses is used by the rational checker. A larger calibrated operational comparison would be needed to decide when the exact certificate cost is justified in practice.
